\subsection{Network Representation}

The public transportation system is represented as a directed
time-expanded graph. All timetable times (arrival/departure event times) are expressed in minutes on the same time axis as the departure-time intervals.

A key modeling choice is (i) the separation of each stop centroid into two
distinct nodes (an origin \emph{centroid-in} and a destination \emph{centroid-out}),
and (ii) the split of each timetable event into an \emph{arrival} and a
\emph{departure} node. This second split prevents ill-defined instantaneous
moves (e.g., transferring or boarding at the exact same time as arrival) and
ensures that all links respect a strict chronological (or lexicographic) order,
so that the resulting time-expanded network is acyclic.

\begin{definition}[Nodes]
The node set consists of three families of nodes:
\begin{itemize}
  \item \textbf{Centroid-in nodes (non time-tagged, indexed by departure interval):}
      for each stop $s\in\mathcal{S}$ and for each departure interval $t\in\mathcal{T}$,
      a node
      \[
      (s,t)^{\mathrm{in}}
      \]
      used as a demand-injection handle for passengers departing from stop $s$
      with desired departure interval $t$. These nodes do not correspond to a physical
      time instant; they are placed \emph{before all event nodes} in the topological
      order (conceptually at $-\infty$). The interval index $t$ is used only to
      (i) select admissible boarding events through the access-window rule, and
      (ii) compute schedule-deviation penalties.
  \item \textbf{Centroid-out nodes (non time-tagged):} for each stop $s\in\mathcal{S}$,
      a node $s^{\mathrm{out}}$ representing the point where passengers arriving at stop
      $s$ exit the network. These nodes are not associated with a physical time; they
      are placed \emph{after all event nodes} in the topological order (conceptually
      at $+\infty$) so that the global graph remains acyclic.
  \item \textbf{Timetable event nodes (trip-specific, split into arrival and departure):}
      for each stop $s$, each scheduled arrival time $\tau^{\mathrm{arr}}$ and departure time
      $\tau^{\mathrm{dep}}$ associated with a given vehicle run (trip), we create two nodes
      \[
      (s,\tau^{\mathrm{arr}})^{\mathrm{arr}}
      \qquad\text{and}\qquad
      (s,\tau^{\mathrm{dep}})^{\mathrm{dep}}.
      \]
      These nodes are \emph{trip-specific}: if two trips serve the same stop at the same
      clock time, they still correspond to different event nodes.
\end{itemize}
\end{definition}

\begin{definition}[Lines and event-to-line mapping]
Let $\mathcal{L}$ denote the set of public transport lines. Each timetable event node
(arrival or departure) is associated with exactly one operated line, captured by the mapping
\[
\lambda:\mathcal{V}^{\mathrm{event}}\rightarrow\mathcal{L},
\qquad v\mapsto \lambda(v),
\]
where $\mathcal{V}^{\mathrm{event}}$ is the set of timetable event nodes.
In practice, $\lambda(v)$ is derived from the timetable: event nodes created from a given
trip inherit the line identifier of that trip. Therefore, for a given trip at stop $s$,
the corresponding arrival and departure nodes share the same line label.
\end{definition}

Centroid-in nodes are not time-tagged: they are placed before all timetable
event nodes (conceptually at $-\infty$) so that access links can connect to
departure events that occur either before or after the desired departure window
without violating acyclicity. Centroid-out nodes are not time-tagged and
are placed after all event nodes (conceptually at $+\infty$), preserving a global
topological order.

